% Generated draft: requires mathematical and source review.
\documentclass[11pt]{article}
\usepackage[margin=25mm]{geometry}
\usepackage{fontspec}
\setmainfont{Latin Modern Roman}
\usepackage{amsmath,amssymb,mathtools,unicode-math}
\setmathfont{Latin Modern Math}
\usepackage{enumitem,needspace,xurl,hyperref}
\setlength{\parindent}{0pt}
\setlength{\parskip}{5pt}
\setlength{\emergencystretch}{3em}
\newcommand{\N}{\mathbb{N}}
\newcommand{\Z}{\mathbb{Z}}
\newcommand{\Q}{\mathbb{Q}}
\newcommand{\R}{\mathbb{R}}
\newcommand{\C}{\mathbb{C}}
\newcommand{\F}{\mathbb{F}}
\newcommand{\A}{\mathbb{A}}
\newcommand{\PP}{\mathbb{P}}
\newcommand{\E}{\mathbb{E}}
\newcommand{\eps}{\varepsilon}
\newcommand{\dd}{\,\mathrm{d}}
\begin{document}
\section*{UnsolvedMath ID: 20002473}
\textbf{Source problem code:} catalogue entry
\subsection{Definitions and mathematical statement}
\textbf{Question (Patrick Bernard).} Can KAM results be proved through regularity of viscosity solutions of the Hamilton--Jacobi cell equation
\[
 H(x,c+Du(x))=\alpha_H(c),\qquad x\in\mathbb T^n?
\]
Here \(T^*\mathbb T^n\cong\mathbb T^n\times\R^n\), \(c\in H^1(\mathbb T^n;\R)\cong\R^n\), and in the Tonelli framework \(\alpha_H(c)\) is the critical value. The source calls existence of KAM invariant tori equivalent to regular solutions. Its regularity, Hamiltonian class, nondegeneracy, frequency and smallness assumptions are unspecified. The aid qualifies that scalar regularity gives an invariant graph, whereas KAM quasiperiodicity requires extra dynamics.

\textbf{Reported graph criterion.} For a \(C^2\) Hamiltonian and periodic \(u\in C^2(\mathbb T^n)\), the closed Lagrangian graph
\[
 \Lambda_{c,u}=\{(x,c+Du(x)):x\in\mathbb T^n\}
\]
is invariant under the Hamiltonian vector field if and only if \(H(x,c+Du(x))\) is constant. This does not alone say that its flow is conjugate to a constant irrational-frequency translation. Such a conjugacy requires, for example, a torus diffeomorphism \(V\) satisfying
\[
 DV(x)\,H_p(x,c+Du(x))=\omega,
\]
with the intended frequency and regularity conditions. A Diophantine frequency satisfies \(|k\cdot\omega|\geq\gamma|k|^{-\tau}\) for all nonzero \(k\in\Z^n\), with specified \(\gamma>0,\tau\geq0\).

\textbf{Reported arbitrary-dynamics realization.} Fix a smooth periodic \(u_0\), \(c\in\R^n\), the closed covector field \(q(x)=c+Du_0(x)\), and any smooth vector field \(b\) on the torus. With the flat fiber norm, define
\[
 H_{b,q}(x,p)=\tfrac12|p-q(x)|^2+
 \langle b(x),p-q(x)\rangle.
\]
The aid reports that this Hamiltonian is uniformly fiberwise convex and Tonelli, that \(\alpha_{H_{b,q}}(c)=0\), and that \(u_0\) is a smooth critical cell solution. Its characteristic dynamics on the graph is exactly
\[
 \dot x=H_p(x,q(x))=b(x).
\]
Thus a smooth scalar cell solution can coexist with arbitrary smooth torus dynamics; no KAM-type quasiperiodicity follows from that regularity alone. This does not refute KAM theorems with their nondegeneracy, perturbative and frequency hypotheses.

\textbf{Reported first-order obstruction.} In a perturbative model \(H_\epsilon(x,p)=h(p)+\epsilon P(x,p)\), with \(\omega=Dh(c)\), expansions \(u_\epsilon=\epsilon u_1+\cdots\), \(\alpha_{H_\epsilon}(c)=h(c)+\epsilon a_1+\cdots\) lead to the cohomological equation
\[
 \omega\cdot Du_1=a_1-P(x,c),\qquad
 a_1=\langle P(\cdot,c)\rangle.
\]
The aid reports a Diophantine analytic strip-loss estimate for this first-order equation but gives no constants, analytic norms, widths or full iteration. It is a linearized calculation, not an implicit proof of a complete nonlinear KAM theorem.

The aid says Hamilton--Jacobi/configuration-space KAM proofs already give an affirmative answer to the broad methodological question, while it found no generic viscosity-to-smooth KAM bootstrap in the literature checked. No titles, dates or search scope are supplied for that assessment. It is retained as a qualified report, without claiming a new absence result or a universal equivalence of smooth cell solutions and KAM dynamics.
\subsection{Short English statement}
The question seeks KAM proofs through Hamilton--Jacobi regularity. The aid distinguishes invariant Lagrangian graphs from quasiperiodic KAM tori: its uniformly convex realization places any smooth torus dynamics on a smooth critical graph. The conjugacy equation and small-divisor estimates remain additional requirements.
\subsection{Sources}
\begin{itemize}
\item Local UnsolvedMath registry, permanent ID \(20002473\), catalogue entry, \texttt{aim-pdes-notes.json}, \texttt{notes[84]}, workshop \emph{New connections between dynamical systems and PDE's}: \url{https://aimath.org/WWN/dynpde/dynpde.pdf}.
\item The supplied verification identifies Problem 10, Patrick Bernard, in the original PDF version dated 15 August 2003. The extraction labels it ``0'' and includes a trailing ``1'' page artifact; these do not alter permanent identity.
\item Supplied research summary on graph invariance, arbitrary-dynamics realization, conjugacy and the cohomological equation, with an unspecified literature assessment of Hamilton--Jacobi/configuration-space KAM proofs. No full authoritative references are supplied. Attempt 1: partial\_result, partially\_solved, rigor ``proved'', low novelty confidence. The machine-generated aid is not peer reviewed and requires independent verification before citation; no PDF visit is claimed.
\end{itemize}
\end{document}
